\section{The actual intrinsic ample cone generates every original numerical class}
\label{sec:native-numerical-ample-cone-generation}

\begin{lemma}[Native projectivity supplies an actual geometric ample line {\cite[\S1.2]{BJ22}}]
\label{lem:native-projectivity-geometric-ample-line}
\lean{EnvelopeProperty.NativeRealNumericalClasses.exists_geometric_ample_line}
\leanok
\uses{def:fc-original-geometric-line-ampleness,lem:fc-original-very-ample-is-ample}
Every native projective structure morphism to the original coefficient
spectrum admits an internally produced actual geometrically ample line.
\end{lemma}
\begin{proof}
Native projectivity gives a genuine finite polynomial projective closed
embedding over the same original coefficient ring. Take the actual
whole pullback of its original $\mathcal O(1)$. The same embedding and
identity line isomorphism prove that this line is genuinely very ample.
Its actual first tensor power proves geometric ampleness.
\end{proof}

\begin{theorem}[Every original numerical class is a difference of intrinsic ample classes {\cite[\S1.2]{BJ22}}]
\label{thm:native-numerical-ample-cone-generation}
\lean{EnvelopeProperty.NativeRealNumericalClasses.exists_ampleCone_difference, EnvelopeProperty.NativeRealNumericalClasses.ampleCone_isReproducing, EnvelopeProperty.NativeRealNumericalClasses.span_ampleCone_eq_top}
\leanok
\uses{lem:native-projectivity-geometric-ample-line,thm:native-numerical-finite-ample-openness,thm:native-numerical-line-family-maps}
On integral projective $X$, every class in the original real numerical
quotient is a difference of two classes of the same intrinsic ample
cone. Consequently that cone is reproducing and its real linear span
is the entire original numerical quotient. No global finite-dimensional
span or numerical ampleness-invariance assumption occurs.
\end{theorem}
\begin{proof}
Produce an actual ample $A$ by the preceding lemma. Represent the given
original $\theta$ in a finite original line family as $F(a)$, using the
actual finite numerical-line representation theorem. Finite-family
ample openness at $[A]$ supplies a positive coefficient radius
$\varepsilon$. Put $S=\sum_i|a_i|$ and
$r=\varepsilon/(S+1)>0$. Since $|a_i|\leq S$, every coefficient of
$ra$ has absolute value less than $\varepsilon$. Hence
$[A]+r\theta$ belongs to the original ample cone. Positive scaling
also puts $r^{-1}([A]+r\theta)$ and $r^{-1}[A]$ in that cone, and their
exact difference is the original $\theta$.
Thus every element belongs to the set difference of the cone with
itself, proving the reproducing property. Its linear span is the whole
space because every such difference belongs to that span. The same
argument covers the zero-dimensional or zero-vector-space cases and
the empty representing family; it uses no artificial nontriviality
assumption on the original numerical quotient.
\end{proof}
